\section{Quartic and cubic canonical realizations of BA dynamics}
\label{sec:AHH-comparison}

The moment-map construction of Adams, Harnad and Previato pulls
matrix Hamiltonians back through bilinear maps; Adams, Harnad and
Hurtubise integrate
both the matrix flows and their canonical lifts using spectral line
bundles and generalized BA functions
\cite{AdamsHarnadPreviato1988,AdamsHarnadHurtubise1990II}.
Here we identify a quartic specialization whose bilinear variables
satisfy precisely our cubic multiwave equations after a fixed change
of time and phase. A complete-fiber BA identity explains why these
bilinear sums reduce to individual endpoint BA functions. This gives
an explicit connection between the two canonical realizations.

\subsection{A quartic canonical parent}

Let $\mathsf F,\mathsf G\in\C^{q\times d}$, with
$\{F_{ih},G_{jk}\}=-i\delta_{ij}\delta_{hk}$ for
$1\leq i,j\leq q$, $1\leq h,k\leq d$, and all other basic brackets
zero. The entries of $\mathsf G$ are independent partners of those
of $\mathsf F$. Set
\begin{equation}\label{eq:AHH-quartic-parent}
 \mathsf M=\mathsf G^{\trans}\mathsf F,\qquad
 \mathcal H_4=\tr(\mathsf D\mathsf M^2),\qquad
 \mathsf C=\mathsf D\mathsf M+\mathsf M\mathsf D.
\end{equation}
Here $\mathsf D$ is the diagonal matrix in
\cref{eq:canonical-pencil}; $\mathcal H_4$ is a selected pullback in
the moment-map framework of \cite{AdamsHarnadPreviato1988}, rather
than a separately distinguished model in that paper.
In scalar modes its full dependence is
\begin{equation}\label{eq:AHH-quartic-components}
 \mathcal H_4=
 \sum_{h=1}^{d}\sum_{k=1}^{d}\sum_{i=1}^{q}\sum_{j=1}^{q}
 D_{hh}G_{ih}F_{ik}G_{jk}F_{jh}.
\end{equation}
Hamilton's equations and their closed bilinear projection are
\begin{equation}\label{eq:AHH-parent-rates}
 i\frac{d\mathsf F}{ds}=\mathsf F\mathsf C,\qquad
 i\frac{d\mathsf G}{ds}=-\mathsf G\mathsf C^{\trans},\qquad
 \widehat{\mathsf X}=-\mathsf G^{\trans}\mathsf F,
 \qquad i\frac{d\widehat{\mathsf X}}{ds}
       =[\mathsf D,\widehat{\mathsf X}^{,2}].
\end{equation}
The diagonal blocks of $\widehat{\mathsf X}$ are constant because
$\mathsf D$ is scalar on each of the three groups. Fix them as in
\cref{eq:canonical-X} and put
\begin{equation}\label{eq:AHH-time-phase-map}
 \mathsf K_{\rm ph}=\diag(-3\mathsf B_L,3\mathsf B_M,0_N),\qquad
 \mathsf X(t)=e^{-it\mathsf K_{\rm ph}}
       \widehat{\mathsf X}(-t)e^{it\mathsf K_{\rm ph}}.
\end{equation}
Then
\begin{equation}\label{eq:AHH-cubic-projection}
 i\dot{\mathsf X}=-[\mathsf D,\mathsf X^2]
                         +[\mathsf K_{\rm ph},\mathsf X],
\end{equation}
whose off-diagonal blocks are exactly
\cref{eq:canonical-matrix-rates} for every positive $L,M,N$.
For $q=d$ every initial $\mathsf X$ admits a factorization, for example
$\mathsf F(0)=\Id_d$, $\mathsf G(0)=-\mathsf X(0)^{\trans}$.
Thus the equation-level projection includes arbitrary canonical
initial modes of the cubic system. The adapted BA factorization is
given in Appendix~\ref{sec:AHH-dictionary}.

\subsection{The two matrix brackets}

Both parent mode systems have constant canonical brackets. Their
relationship is seen by distinguishing the two brackets on
$\mathsf X$. In this subsection write $[f,h]_a=i\{f,h\}_a$ to
suppress the common physical factor $-i$. The bilinear projection
from $\mathsf F,\mathsf G$ induces the linear Lie--Poisson bracket
\begin{equation}\label{eq:AHH-two-brackets}
\begin{aligned}
 [X_{hi},X_{jk}]_1&=\delta_{hk}X_{ji}-\delta_{ij}X_{hk},\\
 [X_{hi},X_{jk}]_0&=(D_{hh}-D_{ii})\delta_{hk}\delta_{ij},
 &&1\leq h,i,j,k\leq d.
\end{aligned}
\end{equation}
The second, constant bracket is exactly
\cref{eq:canonical-brackets} under the block scalings of
\cref{eq:canonical-X}. Its same-group blocks are Casimirs, so their
fixed values specify the canonical cubic mode space. With gradients
defined by $dh=\tr((\nabla h)d\mathsf X)$, the two brackets give
\begin{equation}\label{eq:AHH-Lenard}
 [\mathsf X,\tr(\mathsf D\mathsf X^2)]_1
 =[\mathsf X,\tr(\mathsf X^3)/3]_0
 =[\mathsf D,\mathsf X^2].
\end{equation}
They are compatible: translating $\mathsf X$ by $a\mathsf D$ in the
first bracket gives $P_1(\mathsf X+a\mathsf D)=P_1(\mathsf X)-aP_0$.
This is the argument-shift Poisson pencil already underlying
\cref{sec:canonical}; it explains the change from a quartic pullback
to a cubic canonical Hamiltonian. Indeed, up to an additive constant,
\begin{equation}\label{eq:AHH-cubic-H}
 H=-\frac13\tr(\mathsf X^3)
   +3\sum_{\ell=1}^{L}b_\ell^L Q_\ell^L
   +3\sum_{m=1}^{M}b_m^M Q_m^M.
\end{equation}
The last two terms generate precisely the phase conjugation in
\cref{eq:AHH-time-phase-map}. This compares vector fields on two
Poisson realizations: the bilinear map is many-to-one, and the
canonical spaces have $qd$ and $LM+MN+NL$ pairs, respectively.

\subsection{BA convolution and interaction factorization}

The functional counterpart of the bilinear map is a composition
identity for the Szeg\H{o} kernel
\cite[Proposition 5.2.1]{BenZviBiswas2003}. In scalar local
half-differential frames it is
\begin{equation}\label{eq:AHH-Szego}
 \mathscr S_z(P,Q)=
 \frac{\theta(z+A(P)-A(Q))}{\theta(z)E(P,Q)}.
\end{equation}
This is the theta/prime-form kernel underlying the BA pair functions
used earlier; fixed frames and exponential BA factors restore their
chosen normalization. Let $\Lambda,W$ be meromorphic functions
with poles only at distinct points $P_1,\ldots,P_d$, with $\Lambda$
having a simple pole at each. Choose a complete regular fiber
$\Lambda^{-1}(\alpha)=Q_1+\cdots+Q_d$. In coordinates
$t_h=1/\Lambda$ at $P_h$ and $s_i$ at $Q_i$, write
\begin{equation}\label{eq:AHH-fiber-data}
 \frac{W}{\Lambda-\alpha}
       =y_h+c_ht_h+O(t_h^2),\qquad
 \gamma_i=\frac{W(Q_i)}{\partial_{s_i}\Lambda(Q_i)},
 \quad 1\leq h,i\leq d.
\end{equation}
Assume $W$ has at most simple poles at the $P_h$. Then
\begin{equation}\label{eq:AHH-fiber-convolution}
 \sum_{i=1}^{d}\gamma_i
      \mathscr S_z(P_h,Q_i)\mathscr S_z(Q_i,P_k)
 =\begin{cases}
 (y_h-y_k)\mathscr S_z(P_h,P_k),&h\ne k,\\
 c_h,&h=k.
 \end{cases}
\end{equation}
The contracted half-differentials at $Q_i$ make the formula independent
of the coordinates $s_i$. To see it directly, take residues in $Q$
of $(W/(\Lambda-\alpha))\mathscr S_z(P_h,Q)
\mathscr S_z(Q,P_k)$. For $h\ne k$, the endpoint residues are
$-y_h\mathscr S_z(P_h,P_k)$ and
$y_k\mathscr S_z(P_h,P_k)$. For $h=k$, the kernel product has
expansion $(-t_h^{-2}+O(1))dt_h$, giving residue $-c_h$.
The remaining residues are exactly the displayed sum. This also
proves the cancellation of the extra theta-divisor pole in each
product sum.

For disjoint endpoint lists, assemble
$\mathsf A_{hi}=\mathscr S_z(P_h,Q_i)$,
$\mathsf B_{ik}=\mathscr S_z(Q_i,P_k)$,
$\mathsf K_{hk}=\mathscr S_z(P_h,P_k)$ and
$\mathsf\Gamma=\diag(\gamma_1,\ldots,\gamma_d)$, with
$\mathsf Y_{\rm left},\mathsf Y_{\rm right}$ the corresponding
endpoint diagonals of the $y$ values. The off-diagonal identity is
\begin{equation}\label{eq:AHH-Sylvester}
 \mathsf A\mathsf\Gamma\mathsf B
   =\mathsf Y_{\rm left}\mathsf K-\mathsf K\mathsf Y_{\rm right}.
\end{equation}
This Sylvester displacement relation makes the interaction change
explicit. For compatible vectors $\boldsymbol a,\boldsymbol b$,
\[
 \boldsymbol a^{\trans}
 (\mathsf Y_{\rm left}\mathsf K-\mathsf K\mathsf Y_{\rm right})
 \boldsymbol b
 =\boldsymbol a^{\trans}\mathsf A\mathsf\Gamma\mathsf B\boldsymbol b.
\]
A cubic expression using endpoint pair modes is thereby a quartic
expression in the intermediate evaluations. The weights and the
complete fiber ensure that both sides remain BA functions. In the
AHH realization this is the function-level counterpart of
\cref{eq:AHH-parent-rates}; changing the list of independent modes
also requires their brackets and evolution, supplied here by
\cref{eq:AHH-two-brackets,eq:AHH-cubic-projection}.
